% !TEX root = ../main.tex
\chapter{Design Thinking : brainstorming et prototype}\label{annexe:design-thinking}

% Planches générées par design-thinking/_source (npm run design), copiées dans figures/design-thinking/.

\section{Séance de brainstorming}
La figure~\ref{fig:brainstorming} présente le tableau de la séance de brainstorming du Sprint 0, organisée avec l'équipe de \entreprise. Les idées sont regroupées en dix thèmes autour de l'idée centrale du projet ; elles ont alimenté l'étude de l'existant (chapitre~\ref{ch:contexte}) et le backlog du produit (chapitre~\ref{ch:besoins}).

\insererfigurepaysage{design-thinking/brainstorming.png}{Tableau de brainstorming de \nomprojet}{fig:brainstorming}

\section{Prototype basse fidélité}
La figure~\ref{fig:prototype} présente les maquettes basse fidélité validées avec l'encadrant entreprise avant le développement : le parcours de réservation du patient sur mobile (A1 à A4), l'espace clinique (B1, B2), le tableau de bord et la console d'administration (C1, C2). Les écrans réalisés sont présentés au chapitre~\ref{ch:realisation}.

\insererfigure[0.82]{design-thinking/prototype.png}{Prototype basse fidélité de \nomprojet}{fig:prototype}
